\section{Querschnittliche Konzepte}

Dieser Abschnitt beschreibt Konzepte und Entscheidungen, die nicht einem einzelnen
Service zugeordnet sind, sondern das gesamte System betreffen. Sie bilden die
gemeinsame technische Basis für Sicherheit, Betrieb, Datenhaltung und Qualitätssicherung.

\subsection{Sicherheitskonzept}

Das Sicherheitskonzept basiert auf dem Prinzip der minimalen Rechte: Authentifizierung
und Autorisierung sind konsequent voneinander getrennt, zentral über den
\texttt{identity}-Service koordiniert und in jedem Service durchgesetzt.

\subsubsection{Authentifizierung}

Die Authentifizierung ist vollständig an einen externen OIDC-Provider delegiert.
Das System selbst speichert keine Passwörter -- es vertraut ausschliesslich
signierten JWT-Tokens, deren Gültigkeit am API Gateway geprüft wird.

\begin{itemize}[leftmargin=1.5em]
  \item OIDC / OAuth~2.0: Authentifizierung vollständig extern (z.\,B.\ Keycloak, Azure~AD).
  \item JWT-Access-Token (RS256-signiert, 15~Min.), Refresh-Token (8~h).
  \item API Gateway validiert JWT-Signatur via JWKS-Endpoint; ungültig $\to$ 401.
  \item Silent Refresh: Frontend erneuert Token automatisch.
\end{itemize}

\subsubsection{Autorisierung (RBAC)}

Die Zugriffskontrolle erfolgt rollenbasiert (RBAC). Rollen werden als Claims im
JWT übertragen und von jedem Service eigenständig ausgewertet -- ein zentraler
Autorisierungs-Roundtrip entfällt damit pro Request. Der folgende
Payload-Ausschnitt zeigt die Struktur eines dekodiertes JWT:

\begin{lstlisting}
{
  "sub": "user-uuid-123",
  "roles": ["PARTICIPANT"],
  "locale": "de",
  "tenant_id": "skillbridge",
  "exp": 1700000900
}
\end{lstlisting}

E-Mail ist im OIDC-Token vorhanden, wird aber nie geloggt (PII-Schutz, Kap.~8.6).

Die folgende Tabelle zeigt, welche Aktionen den einzelnen Rollen erlaubt sind:

\begin{tabularx}{\textwidth}{X C{1.8cm} C{1.8cm} C{1.8cm} C{2.2cm}}
  \rowcolor{tableheader}
  \color{white}\textbf{Aktion} & \color{white}\textbf{Admin} & \color{white}\textbf{Staff} & \color{white}\textbf{Trainer} & \color{white}\textbf{Participant} \\
  Kurs erstellen & \yes & \yes & \no & \no \\
  \rowcolor{tablegrey}
  Kurs veröffentlichen & \yes & \no & \no & \no \\
  Kurs archivieren & \yes & \no & \no & \no \\
  \rowcolor{tablegrey}
  Durchführung planen & \yes & \yes & \no & \no \\
  Sich einschreiben & \yes & \yes & \yes & \yes \\
  \rowcolor{tablegrey}
  Eigene Stornierung & \yes & \yes & \yes & \yes \\
  Alle Einschreibungen sehen & \yes & \yes & \no & \no \\
  \rowcolor{tablegrey}
  Nutzer verwalten & \yes & \no & \no & \no \\
\end{tabularx}

\begin{hinweis}[RBAC -- Einschreiberecht für TRAINER]
  Die Berechtigung \emph{„Sich einschreiben"} für die Rolle TRAINER ist bewusst gewählt: Trainer können
  Weiterbildungsangebote ihrer Kollegen als Teilnehmer besuchen (Beispiel: ein Java-Trainer schreibt sich
  in einen Python-Kurs ein). Diese Rollenakkumulation ist im \texttt{identity}-Service durch die
  Multi-Rollen-Unterstützung abgebildet (\texttt{"roles": ["TRAINER", "PARTICIPANT"]} im JWT möglich).
  Falls ein Kursanbieter dies nicht wünscht, kann die Berechtigung per Konfiguration deaktiviert werden.
\end{hinweis}

Die folgende Tabelle beschreibt die eingesetzten Resilience-Muster, die das System
gegenüber transienten Fehlern und Kaskadenausfällen absichern:

\begin{tabularx}{\textwidth}{L{3.5cm} X}
  \rowcolor{tableheader}
  \color{white}\textbf{Muster} & \color{white}\textbf{Beschreibung} \\
  Circuit Breaker & Resilience4j: 5~Fehler in 10\,s $\to$ Open; Half-Open nach 30\,s. \\
  \rowcolor{tablegrey}
  Retry + Exp.\ Backoff & Bis zu 3×; Backoff: 100\,ms / 500\,ms / 2\,s. \\
  Dead-Letter-Queue & Nach 3~Fehlversuchen; Alerting bei DLQ-Tiefe $>$ 0. \\
  \rowcolor{tablegrey}
  Graceful Degradation & Ausfall von \texttt{notification} blockiert nie die Einschreibung. \\
  Idempotente Handler & EventId-Prüfung in \texttt{notification\_log}; Duplikate werden ignoriert. \\
  \rowcolor{tablegrey}
  Timeout & Connect-Timeout 2\,s; Read-Timeout 5\,s. \\
  Bulkhead & Separate Thread-Pools für \texttt{identity}- und \texttt{catalog}-Aufrufe. \\
\end{tabularx}



\textbf{Akzeptiertes Restrisiko: Kein mTLS in v1.0 (TS1)}

In v1.0 läuft die Service-zu-Service-Kommunikation innerhalb des
Kubernetes-Clusters über \textbf{unverschlüsseltes HTTP}.
Ein Angreifer mit Zugang zum internen Netz (z.\,B.\ durch einen
kompromittierten Pod) könnte den Datenverkehr zwischen Services mitlesen.

\textbf{Warum ist das für v1.0 tolerierbar?}
\begin{itemize}[leftmargin=1.5em]
  \item Das interne Kubernetes-Netzwerk ist durch \textbf{Network Policies}
  segmentiert: Jeder Service kommuniziert nur mit explizit erlaubten Peers.
  \item Der Perimeter ist durch TLS am API Gateway gesichert;
  kein externer Traffic erreicht die Services direkt.
  \item Das Betriebsteam hat in v1.0 kein ausreichendes Know-how für
  Sidecar-Proxies und Zertifikatsrotation (Budget- und Zeitrahmen MVP).
\end{itemize}

\textbf{Unter welchen Bedingungen muss nachgebessert werden?}
\begin{itemize}[leftmargin=1.5em]
  \item Vor Go-Live eines zweiten Tenants (Multi-Tenancy v2.0, ADR-008):
  Datentrennung auf Netzwerkebene wird dann zwingend.
  \item Bei einem Security-Audit-Befund, der dieses Muster als
  Critical/High einstuft (Kap.~1.1, Go-Live-Kriterium).
  \item Spätestens mit Einführung von v2.0:
  Linkerd ersetzt das fehlende mTLS vollständig (ADR-007, TS1).
\end{itemize}
Dieses Risiko ist als technische Schuld \textbf{TS1} erfasst und im
Tilgungsplan verankert (Kap.~11).

\subsection{Logging und Observability}

Observability ist als Querschnittsanforderung in allen Services einheitlich
umgesetzt. Strukturierte Logs, verteiltes Tracing und Prometheus-Metriken bilden
zusammen die Grundlage für Fehleranalyse, Alerting und Kapazitätsplanung im Betrieb.

\begin{tabularx}{\textwidth}{L{3cm} X}
  \rowcolor{tableheader}
  \color{white}\textbf{Aspekt} & \color{white}\textbf{Beschreibung} \\
  Strukturiertes Logging & JSON-Format (Logback + SLF4J): \texttt{timestamp, level, service, correlationId, message, exception}. \\
  \rowcolor{tablegrey}
  Correlation-ID & UUID (\texttt{X-Correlation-ID}) wird als HTTP-Header propagiert und in alle Logs geschrieben. \\
  PII-Schutz & E-Mail und Name werden vor dem Logging maskiert (Log-Scrubbing-Middleware). \\
  \rowcolor{tablegrey}
  Metriken & Prometheus (RED-Muster); exportiert via \texttt{/actuator/prometheus}. \\
  Distributed Tracing & OpenTelemetry; Traces in Jaeger; \texttt{traceparent}-Header-Propagation. \\
  \rowcolor{tablegrey}
  Alerting & Error-Rate $>$ 1\,\%, P99-Latenz $>$ 1\,s, Queue-Tiefe $>$ 500, DLQ $>$ 0, DB-Pool $>$ 90\,\%, SSL-Expiry $<$ 30~Tage. \\
\end{tabularx}

\subsection{Datenpersistenz}

Jeder Service verwaltet seine Daten vollständig autonom in einer eigenen
PostgreSQL-Instanz. Kein Service greift direkt auf die Datenbank eines anderen
zu -- Datenaustausch erfolgt ausschliesslich über Events. Dieses Muster
garantiert lose Kopplung und ermöglicht unabhängige Schema-Migrationen.

\begin{tabularx}{\textwidth}{L{3.5cm} X}
  \rowcolor{tableheader}
  \color{white}\textbf{Konzept} & \color{white}\textbf{Beschreibung} \\
  Database per Service & Eigene PostgreSQL-DB je BC; kein Cross-Service-Zugriff. \\
  \rowcolor{tablegrey}
  Optimistic Locking & \texttt{version}-Feld im Aggregate; Konflikt $\to$ 409~Conflict + Client-Retry. \\
  Transactional Outbox & Events atomar mit Domänenänderung; Relay-Prozess publiziert asynchron (vgl.\ ADR-009). \\
  \rowcolor{tablegrey}
  Schema-Migration & Flyway; versionierte SQL-Skripte; idempotent und rückwärtskompatibel. \\
  Backups & Täglich; PITR~7~Tage; wöchentliche Restore-Tests. \\
\end{tabularx}

\subsection{API-Design-Richtlinien}

Alle Services folgen denselben API-Design-Richtlinien, um eine konsistente
Schnittstelle für Clients und Service-zu-Service-Kommunikation sicherzustellen.
Grundlage ist ein \textbf{Contract-First-Ansatz} (ADR-011): Die
OpenAPI-Spezifikation (\texttt{openapi/}\allowbreak \texttt{eduplanner-}\allowbreak \texttt{api.yaml}) ist die
\emph{einzige Quelle der Wahrheit} -- Interfaces und Server-Stubs werden
daraus generiert, nicht umgekehrt.
Die vollständige Spezifikation liegt in Anhang~D.

\subsubsection{Contract-First-Workflow}

Im Contract-First-Ansatz wird die Reihenfolge gegenüber Code-First
umgekehrt: Zuerst entsteht der Vertrag, dann der Code.

\begin{enumerate}[leftmargin=2em]
  \item \textbf{Spec schreiben:} Die Datei
  \texttt{openapi/}\allowbreak \texttt{eduplanner-}\allowbreak \texttt{api.yaml} wird manuell gepflegt und
  in Git versioniert. Sie ist der einzige Ort, an dem API-Änderungen
  initiiert werden.
  \item \textbf{Stubs generieren:} Der \texttt{openapi-}\allowbreak \texttt{generator-}\allowbreak \texttt{maven-}\allowbreak \texttt{plugin}
  erzeugt beim Build automatisch Spring-Boot-Interfaces (Server-Stubs)
  aus der Spec:
\begin{lstlisting}
<!-- pom.xml (Auszug) -->
<plugin>
  <groupId>org.openapitools</groupId>
  <artifactId>openapi-generator-maven-plugin</artifactId>
  <configuration>
    <inputSpec>
      ${project.basedir}/../../openapi/eduplanner-api.yaml
    </inputSpec>
    <generatorName>spring</generatorName>
    <configOptions>
      <interfaceOnly>true</interfaceOnly>
      <useSpringBoot3>true</useSpringBoot3>
      <useTags>true</useTags>
    </configOptions>
  </configuration>
</plugin>
\end{lstlisting}
  \item \textbf{Implementieren:} Das Dev-Team implementiert die
  generierten Interfaces (\texttt{CoursesApi}, \texttt{RegistrationsApi}
  usw.). Der Compiler erzwingt Konformität -- eine Abweichung vom
  Vertrag führt zu einem Build-Fehler.
  \item \textbf{Swagger~UI bereitstellen:} Jeder Service stellt die
  interaktive Swagger~UI unter
  \texttt{/swagger-ui/index.html} bereit (via
  \texttt{springdoc-}\allowbreak \texttt{openapi-}\allowbreak \texttt{starter-}\allowbreak \texttt{webmvc-}\allowbreak \texttt{ui}).
  Die zugrundeliegende YAML-Datei ist unter
  \texttt{/v3/api-docs.yaml} abrufbar.
  \item \textbf{CI-Validierung:} Die Pipeline validiert die Spec vor
  jedem Merge mit \texttt{vacuum} (Linting) und
  \texttt{openapi-generator} (Generierbarkeit). Ein ungültiger
  Vertrag blockiert den Build.
\end{enumerate}

\begin{reviewbox}[Contract-First -- Vorteile für EduPlanner]
  \textbf{Frühe Einigung:} Frontend- und Backend-Team können parallel
  entwickeln, sobald der Vertrag feststeht -- kein Warten auf
  funktionierende Implementierungen.\\[4pt]
  \textbf{Konsistenz:} Generierte Stubs und DTOs sind immer synchron
  mit der Spec. Manuelle Abweichungen sind nicht möglich.\\[4pt]
  \textbf{Contract Testing (TS2):} Pact-Tests validieren, dass
  Consumer und Provider denselben Vertrag einhalten
  (Tilgung Sprint~8, Kap.~11.2).
\end{reviewbox}

\subsubsection{API-Design-Regeln}

\begin{itemize}[leftmargin=1.5em]
  \item \textbf{RESTful:} Ressourcenorientierte URLs, korrekte HTTP-Verben
  (GET\,/\,POST\,/\,PUT\,/\,DELETE).
  \item \textbf{Versionierung:} \texttt{/api/v1/...} --
  Breaking Changes führen zu einer neuen Major-Version
  (\texttt{/api/v2/...}). Additive Änderungen sind rückwärtskompatibel.
  \item \textbf{Swagger~UI:} Erreichbar unter
  \texttt{/swagger-ui/index.html} pro Service. Zugriffsschutz in
  Produktion via API-Gateway-Policy (nur internes Netz).
  \item \textbf{Fehlerantworten:} RFC~7807 Problem Details
  (\texttt{application/problem+json}) mit \texttt{type},
  \texttt{title}, \texttt{status}, \texttt{detail}, \texttt{instance}
  und optionalem \texttt{errors}-Array für Validierungsfehler.
  \item \textbf{Pagination:} Cursor-basiert
  (\texttt{?cursor=\ldots\&limit=20}, max.\ 100).
  Kein Offset-Paging (inkonsistent bei hoher Schreiblast).
  \item \textbf{Idempotency-Key:} Optionaler \texttt{Idempotency-Key}-Header
  (UUID) bei allen POST-Endpunkten -- verhindert Doppeleinschreibungen
  bei Netzwerkfehlern (vgl.\ Kap.~5.4).
  \item \textbf{Authentifizierung:} RS256-signierter JWT Bearer Token;
  Validierung am API Gateway; Services vertrauen dem weitergeleiteten
  Header (Defense in Depth via lokales RBAC).
\end{itemize}

\begin{hinweis}[Beispiel: RFC~7807 Fehlerantwort]
\begin{lstlisting}
{
  "type":     "https://eduplanner.skillbridge.ch/problems/validation-error",
  "title":    "Validation Error",
  "status":   400,
  "detail":   "Das Feld 'title' darf nicht leer sein.",
  "instance": "/api/v1/courses",
  "errors":   [{ "field": "title", "message": "Darf nicht leer sein." }]
}
\end{lstlisting}
\end{hinweis}

\subsection{Teststrategie (Ebene 1--3)}

Die Teststrategie folgt der Testpyramide: Schnelle, isolierte Unit-Tests bilden
die Basis; Integration- und Contract-Tests sichern das Zusammenspiel der
Komponenten; E2E- und Performance-Tests validieren das Systemverhalten unter
realistischen Bedingungen. Security- und Chaos-Tests ergänzen die Strategie
um Robustheit und Angriffssicherheit.

\begin{tabularx}{\textwidth}{L{2.2cm} L{2.5cm} X L{3.0cm}}
  \rowcolor{tableheader}
  \color{white}\textbf{Testtyp} & \color{white}\textbf{Umfang} & \color{white}\textbf{Ziel} & \color{white}\textbf{Tool / Tooling} \\
  Unit Tests & Einzelne Klassen / Aggregate & Functional Correctness; Business-Logik-Abdeckung $>$ 80\,\% & JUnit~5, Mockito. \\
  \rowcolor{tablegrey}
  Integration Tests & Slice (Controller $\to$ Repo) & Prüfung DB-Interaktion, JPA-Mappings, JSON-Serialisierung & Testcontainers (PostgreSQL). \\
  Contract Tests & Service-Schnittstellen (CDC) & API-Kompatibilität zwischen catalog/sched/reg (Tilgung TS2) & Pact (geplant ab Sprint~8). \\
  \rowcolor{tablegrey}
  E2E Tests & Krit.\ Journeys & Staging-Deploy bricht bei Fehler & Cypress; Einschreibung, Stornierung, Veröffentlichung. \\
  Performance & P95 $<$ 500\,ms & Staging-Alert bei Überschreitung & k6, 1\,000 concurrent users, wöchentlich. \\
  \rowcolor{tablegrey}
  Chaos Engineering & Wöchentlich & Ergebnisse im Wiki & Pod-Kills, Network-Delays (Chaos Mesh). \\
  Security Tests & Vor Release & Release blockiert bei Critical/High & OWASP ZAP, Trivy, Snyk. \\
\end{tabularx}